\documentclass[10pt]{article}
\usepackage[margin=0.8in]{geometry}
\usepackage{lmodern}
\usepackage{microtype}
\usepackage[hidelinks]{hyperref}
\setlength{\parindent}{0pt}
\setlength{\parskip}{0.48em}

\begin{document}

\begin{flushright}
12 August 2026\\
Roberto Fern\'andez-Barrios\\
Faculty of Engineering, University of Deusto\\
Avda. de las Universidades 24, 48007 Bilbao, Spain\\
\href{mailto:roberto.fernandez.b@deusto.es}{roberto.fernandez.b@deusto.es}
\end{flushright}

Editors\\
\emph{npj Quantum Information}

\textbf{Re: Article submission to the Collection ``Quantum machine learning:
understanding capabilities, limitations, and perspectives for quantum
advantage''}

Dear Editors,

We submit ``Conditional validity of quantum event classifiers under collider
systematics and quantum estimation uncertainty'' for consideration as an
Article in the Collection named above. Rather than presupposing quantum
advantage, the manuscript asks a prior question for deployable QML: under what
observable information can a performance or scientific-inference claim be
supported when both the scientific environment and the quantum estimator are
uncertain?

The work makes a connection that is not captured by nominal QML benchmarking.
Collider systematics determine whether the scientific claim is identifiable;
finite-shot and noisy quantum-kernel evaluation additionally randomizes the
implemented Gram matrix, refit, calibration, and operating point. We integrate
these layers in one fail-closed framework and distinguish claims about the
realized deployment from claims anchored to an ideal quantum kernel. This
allows the paper to say exactly what quantum measurement changes without
attributing ordinary stochasticity exclusively to quantum systems.

The results directly address the Collection's emphasis on capabilities and
limitations. Matched classical controls remove every apparent
quantum-specific performance and sensing effect, so we claim no quantum
advantage. What remains quantum-specific is an experimentally quantified,
measurement-induced deployment uncertainty whose effects propagate through
the full decision pipeline. The collider setting is essential as a stringent
scientific deployment environment: it supplies physics weights, invisible
yield shifts, control regions, finite templates, and downstream
signal-strength inference. Evidence spans disjoint simulated worlds, CMS open
data, and a deliberately micro-scale IBM QPU demonstration, with all scope
limits stated explicitly.

The manuscript combines formal identifiability and anytime-valid per-claim
guarantees with registered falsifiers, immutable manifests, public code and
data artifacts, and a complete correction trail. Several registered
falsifiers fired; each adverse result was retained and used to narrow the
claims. We believe this combination of rigorous negative controls, quantum
deployment semantics, and physics-level validation will be useful to the
journal's interdisciplinary QML readership.

For transparency, the same author group has released a distinct manuscript,
``Sharp Target-Domain Certificates for Quantum-Kernel Advantage under
Distribution Shift'' (Zenodo DOI 10.5281/zenodo.21776862). That work studies
finite-target partial identification of predictive advantage against a
prespecified classical-kernel family on unrelated cybersecurity and tabular
datasets. It shares no datasets, experiments, theorems, estimands, codebase, or
reported results with the present collider-inference study. We provide it as
related material so the Editors can assess the distinction directly.

This manuscript is not under consideration elsewhere. All authors have
approved the manuscript and its submission. The authors declare no financial
or non-financial competing interests.

Sincerely,

Roberto Fern\'andez-Barrios\\
Corresponding author, on behalf of all authors

\end{document}
